\section{From the integral connecting cocycle to the ramified character condition}

% BEGIN K81C1
\paragraph{K81C1. A torsion class on the truncated integral period ring.}
Let $p$ be odd, $K_0$ the maximal unramified subfield of $K$, and
$F=K_0(\mu_p)\supset K$. Use the level-one data on the original tied
period rings. Denote the cyclotomic character by $\chi$, the subgroup
$G_F$ by $G_1$, and $G_{K_0(\mu_{p^\infty})}$ by $H$. Fix arbitrary
$e\ge0$ and put $M=e+1$, $A=A_{\inf}$, $A_M=A/(\xi^M)$ and
$B_M=B_{\rm dR}^+/\mathrm{Fil}^{M}$. Then $A_M[1/p]=B_M$ on
these rings: completion in the principal ideal $(\xi)$ induces an
isomorphism on each finite quotient, and $\xi$ generates the kernel
of $\theta$ after inverting $p$. Before inverting $p$ the quotient
embeds, since $p$ is regular modulo $\xi$ and hence modulo its
powers by induction. The map carries $1$ to $1$ and is equivariant.

Take a tested integral short exact sequence $Y\to\widetilde Y\to Z$
from $\mathrm{TBat}$ with trivial quotient line and let $y$ be the
given integral lift. Write $f(g)$ for the unique lattice vector
mapping to $gy-y$. Injectivity gives $f(gh)=f(g)+gf(h)$ and a
linear left inverse to the embedding shows continuity. Over $F$ the
middle term is crystalline by hypothesis. Strictness at filtration
zero (J80H1, or the strict filtered de Rham comparison) gives a
$G_1$-fixed vector in $\widetilde Y\otimes B_{\rm dR}^+$ above
$\bar y\otimes1$. Subtract it from $y\otimes1$. The result lies in
$Y\otimes B_{\rm dR}^+$ by exactness of the underlying vector
spaces, and has coboundary $f\otimes1$. After reduction to $B_M$,
clear the finitely many $p$-denominators of its coordinates in a
$\mathbb Z_p$-basis of the kernel lattice $T$. Thus $f\otimes1$ is
a torsion cocycle in $T\otimes A_M$. Its reductions modulo open
lattice neighborhoods are continuous by continuity of $f$.

Use the integral Wach torsion bound W81T6 at $(a,b)=(0,M)$ for the
$\mathrm{HypW}$ of this actual kernel. In the notation of the
original $\mathrm{cLL}$ put
\[
 N=\mathrm{cLL}\bigl(p,m,\mathrm{vdegMs}(m),J(Y),0,e+1\bigr),
 \qquad m=2+ n_0(p,M_0+h_0)+1.
\]
It yields one $X\in T\otimes A_M$ with
$(g-1)X=p^N f(g)\otimes1$ for every $g\in G_1$. The subsequent
argument uses this identity at the stated $N$.
% END K81C1

% BEGIN K81C2
\paragraph{K81C2. Integral contraction and rational lifting.}
Let $\phi:T\to V/T_e$ be the tested equivariant map, where
$V$ is the coefficient line with character $\chi^{-e}$ and $T_e$
its prescribed lattice. Lift $\phi$ linearly to $\Phi:T\to V$
on a free integral basis. Relative to a $\mathbb Z_p$-basis of
$T_e$ every coordinate $\psi$ of $\Phi$ is $\mathbb Z_p$-linear,
bounded on $T$ and equivariant modulo $\mathbb Z_p$. Contract the
finite tensor $X=\sum t_i\otimes x_i$ with $\psi$. This gives
$b=\sum\psi(t_i)x_i\in B_M$. For $c(g)=\psi(f(g))$ and
$D_g b=\chi(g)^{-e}g b-b$, equivariance modulo $\mathbb Z_p$
gives directly
\[
                 p^N c(g)-D_g b\in A_M\qquad(g\in G_1).
\]
Indeed the error is $\sum(\psi(gt_i)-\chi(g)^{-e}\psi(t_i))gx_i$;
the arbitrary common denominator of $\psi$ costs no power of $p$.
Moreover $c$ modulo $\mathbb Z_p$ is an actual continuous cocycle.

Put $h=1+v_p(e+1)$. By integral Tate duality
$H^2(F,\mathbb Z_p(-e))$ is dual to the invariants of
$\mathbb Q_p/\mathbb Z_p(e+1)$. An element $\gamma\in G_1$
with $\chi(\gamma)=1+p$ acts on these by $(1+p)^{e+1}$, whose
difference from one has valuation $h$. Such a $\gamma$ exists
because $K_0$ is unramified and the cyclotomic tower has its full
procyclic group above level one. Hence $p^h$ kills the obstruction
to lifting the class of $c$. The exact sequence for
$\mathbb Z_p(-e)\subset\mathbb Q_p(-e)$ gives a continuous rational
cocycle $z$ with $z(g)-p^h c(g)\in\mathbb Z_p$ on $G_1$:
a discrepancy by a torsion coboundary is removed by subtracting
the coboundary of any rational lift of its value. The rational
cocycle $Z=p^N z$ consequently satisfies
$Z(g)-D_g(p^hb)\in A_M$ on $G_1$.
% END K81C2

% BEGIN K81C3
\paragraph{K81C3. The coefficient of the logarithmic cocycle and its bound.}
Write $C=\mathbb C_p$, $\delta=\zeta_p-1$, and $t=\log[\epsilon]$.
For a rational $\chi^{-e}$-cocycle $Z$ there is a unique
$\alpha\in F$ such that, in $B_M$, for some $b_0$,
\[
 Z(g)=\log\chi(g)\,\alpha t^e+D_g b_0 .
\]
We give the graded calculation and its integral estimate. For a
continuous unit-ball-valued cocycle on $H$ acting on $C$, multiply
by $p$ to obtain the coboundary of an element of the unit ball.
Indeed in a finite Galois quotient with coset representatives $s$
pick integral $a$ of trace $p$ by A79L1. The negative weighted sum
$-\sum_s s(a)c(s)$ has norm at most the bound on $c$ and its
coboundary differs from $pc$ by at most the bound on the small
open kernel, by permuting the cosets in the cocycle identity.
Continuity makes the latter bound arbitrarily small. Apply this
again to the residual cocycle divided by $p$, choosing successive
errors at most $|p|2^{-j}$; the successive increments have norm
at most $2^{1-j}$ and converge to a primitive of $pc$.

For a unit-ball cocycle of $C(j)$ on $G_1$ the twist is trivial
on $H$. Subtract this integral coboundary from $p$ times the
cocycle. The result factors through $G_1/H$, and its value on
$\gamma$ is in the unit ball of $C^H=\widehat K_\infty$.
Use A79L3's integral splitting into $F$ and the complementary
summand. On the complement $\chi(\gamma)^j\gamma-1$ has inverse
of norm at most $|\delta|^{-1}$. On $F$, if $j\ne0$, divide by
$(1+p)^j-1$, of valuation $1+v_p(j)$. This gives a primitive of
the original cocycle of norm at most $p^{2+v_p(j)}$; equality of
cocycles on $H$ and $\gamma$ implies equality by density.
There are no invariant vectors for $j\ne0$, by the same splitting.
If $j=0$ the residual $F$-coordinate is instead $p\alpha\log(1+p)$.
Thus the unique logarithmic coefficient of a bounded cocycle has
norm at most $p^2$. Uniqueness follows on projecting a putative
logarithmic coboundary to $F$: its primitive is fixed by $H$ and
normalized trace kills $\gamma-1$ there. Compactness permits
scaling any continuous cocycle to the unit ball.

Filter $B_M$ by powers of $t$. At step $i$ the degree-$i$ error
is a $C(i-e)$-cocycle. At $i<e$ lift a primitive via $\theta$
and subtract its coboundary times $t^i$; at $i=e$ the preceding
weight-zero calculation gives the displayed formula. The quotient
topology is the $p$-adic lattice topology on $A_M[1/p]$, and
its graded maps to $C$ are continuous. The Galois action on
these quotients is continuous: on $A$ this follows on finitely
many Witt coordinates modulo $(p^k,\xi^M)$ from continuity on
the tilt in its valuation topology, and extends by inverting $p$.
The same graded descent for the difference of two such expressions
proves uniqueness: the earlier degrees have zero invariants, and
a logarithmic coboundary in the last degree has coefficient zero.

Suppose also that $Z-D_g b$ is integral in $A_M$. The difference
of the two expressions says $\log\chi(g)\alpha t^e$ is an integral
error plus a coboundary with primitive in degree at least zero.
At degree $i<e$ write the primitive as $t^i y$. In the coordinate
relative to $\xi^i$ the residue is $\delta^i\theta(y)$, since
$t/\xi$ has residue $\delta$ by $\mu=\xi([\epsilon^{1/p}]-1)$
and $t/\mu\equiv1\pmod{\mathrm{Fil}^1}$. The integral error in
this coordinate is in $\mathcal O_C$ because
$A\cap\mathrm{Fil}^i=\xi^iA$. Thus the twisted coboundary of
$\delta^i\theta(y)$ is integral. Its unique primitive has norm
at most $p^{2+v_p(e-i)}$. Multiply the entire relation by this
power of $p$, lift the now integral residue by $\theta:A\to\mathcal O_C$,
and subtract $\xi^i$ times the lift from the primitive. Its degree
increases by one with the error still integral. The total power
used in degrees below $e$ is $n=2e+v_p(e!)$. In degree $e$ the
weight-zero logarithmic coefficient of the integral error is
$\delta^e p^n\alpha$, up to sign. With $|\delta|=p^{-1/(p-1)}$
we obtain
\[
                      |\alpha|\le p^{2e+v_p(e!)+2+e/(p-1)}.
\]
% END K81C3

% BEGIN K81C4
\paragraph{K81C4. An integral cocycle with a prescribed power-sum coefficient.}
Write $u=1+p$, $U=u^{-e}$, $E=[\epsilon^{1/p}]$, $\pi=E-1$,
$\mu=[\epsilon]-1=\xi\pi$, and let $\sigma$ be arithmetic
Frobenius on $K_0$. Choose a finite list of pairs $(c_i,m_i)$ with
$c_i\in\mathcal O_{K_0}$ and $m_i$ positive and prime to $p$.
Suppose $d_0\in\mathcal O_{K_0}$ solves
$d_0-\sigma d_0=\sigma(\sum_i c_i)$. For $j\ge1$ define
$d_j=(1-p^j\sigma)^{-1}\sigma(\sum_i c_i m_i^j)$;
the inverse is a convergent geometric series. There exists an
integral continuous $\mathbb Z_p(-e)$-cocycle with logarithmic
coefficient
\[
     \frac{\sum_i c_i m_i^e\zeta_p^{m_i}+p^e d_e}
                       {\log(u)\,e!\,p^e} .
\]
Here and below use the same $p$-power roots for $t$ and $E$.
We prove the integral assertion, including the period convergence.
Embed $\mathcal O_{K_0}$ in $A$ by the constant unramified Witt
vectors. Put $Q=\sum_i c_i E^{m_i}$. There is $y\in A$ with
$(1-\varphi)y=Q$ given by
\[
 y=Q+d_0+
     \sum_{n\ge0}\sum_i\sigma^{n+1}(c_i)
                         \bigl((1+\mu)^{m_i p^n}-1\bigr).
\]
Indeed $\varphi E=1+\mu$, the constant terms give the equation
for $d_0$, and the rest telescopes. The series converges in the
$(p,\mu)$-adic topology: for every fixed degree $k>0$ the
coefficients $\binom{m_i p^n}{k}$ tend to zero $p$-adically,
while high degrees are divisible by a high power of $\mu$.
The ring of integral Witt vectors of the complete valuation ring
of the tilt is complete in this topology (modulo each power of
$p$ use its finitely many Witt coordinates and the valuation
topology). The map to $B_M$ is continuous for this topology since
$\mu^M=0$ there. In $B_M$ one has $E=\zeta_p\exp(t/p)$:
both sides have the same residue, and taking their $p$-th powers
and using the finite logarithm gives equality. The Taylor expansion
of $y$ is $\sum_{j=0}^e w_jt^j$, with
\[
 w_j=\frac{\sum_i c_i m_i^j\zeta_p^{m_i}+p^j d_j}{j!p^j}.
\]
In fact the tail's $j$-th coefficient for $j>0$ is the geometric
series defining $d_j/j!$. This also verifies the expansion when
$e=0$.

We next construct $b\in B_M$ for which $h b-b\in\mathbb Z_p$
for $h\in H$ and $y+U\gamma b-b\in\mathbb Z_p$. To do so take
integral sequences $P_k,Q_k$ with $Q_0=Q$, $P_k$ fixed by $H$,
and in the ring with $\mu$ inverted the identity
\[
  (U\gamma-1)\frac{P_k}{\mu^{k+1}}
       =\frac{Q_k}{\mu^k}-p\frac{Q_{k+1}}{\mu^{k+1}}.
\]
Here is a coefficient construction. Work on one monomial $E^m$,
using $X=\mu$ and $v=((1+X)^u-1)/X$, which is a formal unit.
If $Q_k=E^m f(X)$ set $P_k=E^m P(X)$, $a=k+1$. It suffices to
solve
\[
 U(1+X)^m P(Xv)-v^aP(X)\equiv v^a X f(X)\pmod p .
\]
On $X^i$ the left side has zero coefficient in degrees $\le i$
modulo $p$, and coefficient $m$ in degree $i+1$, since
$v\equiv1\bmod(p,X^2)$ for odd $p$. Successive coefficients
therefore determine $P$ modulo $p$. Lift integrally and divide the
error by $p v^a$, with the sign chosen in the displayed recursion,
to define the next $f$. Start with $f=1$ and evaluate the series
at $\mu$ using completeness. This uses
$\gamma(E)=E(1+\mu)$ and $\gamma(\mu)=\mu v(\mu)$.
Sum with coefficients $c_i$, all fixed by $H$ and $\gamma$.

Choose $\beta_k,\eta_k\in A$ satisfying
\[
 (\varphi-\xi^{k+1})\beta_k=P_k,\qquad
 (\varphi-\xi^k)\eta_k=Q_k,\qquad \eta_0=-y .
\]
For surjectivity, modulo $p$ the polynomial $x^p-a x-b$ with
integral coefficients has an integral root in the algebraically
closed tilt. Subtract the lift, divide the error by $p$, and repeat
by $p$-adic completeness. In $W(C^\flat)$ the element $\pi$ is
invertible and $\varphi\pi=\mu$. The last recursion then says that
\[
 (U\gamma-1)\frac{\beta_k}{\pi^{k+1}}-
       \frac{\eta_k}{\pi^k}+p\frac{\eta_{k+1}}{\pi^{k+1}}
                  \quad\hbox{is fixed by }\varphi .
\]
It belongs to $\mathbb Z_p$ since each Frobenius-fixed Witt coordinate
is in $\mathbb F_p$. Likewise $(h-1)\beta_k/\pi^{k+1}$ is in
$\mathbb Z_p$. These equalities transport to $B_M$ by clearing
powers of $\pi$ and their Galois translates in $A$: all have nonzero
residue and become units in $B_M$.

Put $b=\sum_{k\ge0}p^k\beta_k/\pi^{k+1}$ in $B_M$. This
series converges in its $p$-adic lattice topology. Indeed
$\xi=\pi^{p-1}+p(1+\pi w)$ with $w\in A$, by the binomial
formula. Modulo $(\pi^a,\xi^M)$, $1+\pi w$ is a unit, and for
$n=\lceil a/(p-1)\rceil$ the binomial expansion of
$(\xi-\pi^{p-1})^{n+M-1}$ vanishes. Thus
$p^{n+M-1}\pi^{-a}\in A_M$ inside $B_M$. Since $p>2$ this gives
$p^k\pi^{-k-1}A_M\to0$, and also $p^k\eta_k/\pi^k\to0$.
The ring $A_M$ is $p$-adically complete and separated: the
successive $\xi$-graded modules are $\mathcal O_C$, and all
successive quotients are $p$-torsion free, so that the same finite
filtration is exact modulo $p^n$ and under the surjective inverse
limits. The telescoping identity now proves the assertions for $b$;
$\mathbb Z_p$ is compact, hence closed in $B_M$.

For each $j<e$ add $w_jt^j/(u^{j-e}-1)$ to $b$. These terms
are $H$-fixed, and their twisted differences at $\gamma$ are
exactly the omitted lower terms. For the resulting $b'$ set
\[
 z(g)=D_g b'+\log\chi(g)\,(w_e/\log u)t^e\quad\text{in }B_M .
\]
It is integral scalar on $H$ and $\gamma$, and it is a continuous
cocycle on $G_1$ with values initially in $B_M(-e)$. Since the
scalar lattice $\mathbb Z_p$ is closed and invariant under the
twist, the cocycle identity and density show its values lie there
on all of $G_1$. The subspace topology on this lattice is its
usual one, again by compactness and separation. This proves the
claimed integral formula.
% END K81C4

% BEGIN K81C5
\paragraph{K81C5. The full integral image at level one.}
Set $\beta=-1$ for $e=0$ and $\beta=-e-v_p((e-1)!)$ for $e>0$.
Every element of $p^\beta\mathcal O_F$ is the logarithmic
coefficient of an integral scalar cocycle. To see this write
$x=\sum_{a=1}^{p-1}b_a\zeta_p^a$ in the integral basis over
$\mathcal O_{K_0}$ obtained by multiplying the power basis by
$\zeta_p$. The power basis is integral since $\zeta_p-1$ is
an Eisenstein uniformizer over $K_0$. Solve
$(1-p^{e+1}\sigma)d=\sigma(\sum_a b_a)$ geometrically and set
$\lambda_a=b_a+p^e d$. Then
\[
     d-p^e\sigma d=\sigma\!\left(\sum_a\lambda_a\right),
     \qquad \sum_a\lambda_a\zeta_p^a+p^ed=x,
\]
using $\sum_{a=1}^{p-1}\zeta_p^a=-1$. For $e=0$ use the pairs
$(\lambda_a,a)$ and $d_0=d$ in K81C4. The output is $x/\log u$,
which ranges over $p^{-1}\mathcal O_F$. For $e>0$ use the two
pairs $(-\lambda_a/a^e,a)$ and $(\lambda_a/a^e,ua)$ for each
$a$. Their zeroth moment vanishes, so $d_0=0$ works. Their
$e$-th moment gives $d_e=(u^e-1)d$, and $\zeta_p^{ua}=\zeta_p^a$.
Thus the output is $(u^e-1)x/(\log(u)e!p^e)$. Its scalar
multiplier has valuation $\beta$ because
$v_p(u^e-1)=1+v_p(e)$ and $v_p(\log u)=1$ for odd $p$.

Let $a=e+3+v_p(e)+\lfloor e/(p-1)\rfloor$, with $v_p(0)=0$
in this expression. The bound of K81C3 gives
$p^a\alpha\in p^\beta\mathcal O_F$. For $e>0$ compare exponents
by $v_p(e!)=v_p(e)+v_p((e-1)!)$ and
$e/(p-1)\le 1+\lfloor e/(p-1)\rfloor$; at $e=0$ use $a=3$.
Subtract from $p^aZ$ the integral cocycle just constructed with
coefficient $p^a\alpha$. This gives a rational cocycle $z'$ of
coefficient zero satisfying
\[
 z'(g)\equiv p^{a+N+h}c(g)\pmod{\mathbb Z_p}
                    \qquad (g\in G_1).
\]
% END K81C5

% BEGIN K81C6
\paragraph{K81C6. The zero coefficient and the exact restriction index.}
As a $\mathbb Q_p$-vector space the continuous rational cocycles
$Z^1(F,\mathbb Q_p(-e))$ have dimension $[F:\mathbb Q_p]+1$.
Indeed rational local Tate duality gives $H^2=0$ since $e+1>0$
and $\chi(G_1)$ is infinite. The local Euler formula gives
$\dim H^1=[F:\mathbb Q_p]+\dim H^0$, and the dimension of the
coboundaries is $1-\dim H^0$. The map to $F$ taking the logarithmic
coefficient is linear by K81C3 and surjective by K81C5 and
rational scaling. Its kernel on cocycles therefore has dimension
one. If $e>0$ the nonzero cocycle $g\mapsto\chi(g)^{-e}-1$ has
coefficient zero, with primitive $1$. Thus each element of the
kernel is a rational coboundary. If $e=0$ the nonzero unramified
degree homomorphism has coefficient zero: in the completed unramified
Witt ring the difference of Frobenius and identity is surjective
by successive Artin--Schreier lifting on the residue algebraic
closure. Choose a primitive for its value at residue Frobenius
(using the residue-degree power of Frobenius). It belongs to
$B_{\rm cris}^+\subset B_{\rm dR}^+$; density gives the primitive
equation on the whole unramified quotient, as in S80J4. The
kernel consequently consists of unramified homomorphisms. In both
cases $z'$ from K81C5 becomes a coboundary over $B_{\rm cris}$.

Run K81C2--K81C5 in every coordinate of $T_e$. Combining the
rational cocycles $z'$ and their crystalline primitives in the
chosen basis gives a continuous $V$-cocycle that is finite over
$F$ and reduces to $p^{a+N+h}\phi(f)$ there. Corestrict its
class along $F/K$. Corestriction is a finite sum on cosets,
commutes with reduction and with the map to crystalline period
cohomology, and sends a coboundary to a coboundary. Hence the
corestricted rational class is finite over $K$. On the torsion
class restricted from $K$ its reduction is $[F:K]$ times the
original class, by restriction followed by corestriction. Write
$[F:K]=p^{g_K}d$ with $d$ a $p$-adic unit and divide the rational
finite lift by $d$. Equality of the torsion cohomology classes
allows an additional torsion coboundary, precisely the coboundaries
included in the cocycle definition of $H_f^1(K,V/T_e)$. We obtain
\[
        p^{N+a+h+g_K}\phi(f)\in H_f^1(K,V/T_e).
\]
Here $g_K=v_p([F:\mathbb Q_p])-v_p([K:\mathbb Q_p])$, and
$N+a+h+g_K=\mathrm{cTBY}(Y,M_0,h_0,e)$, since the two
corrections at level one are
$e+2+v_p(e)+\lfloor e/(p-1)\rfloor+1$ and $1+v_p(e+1)$.
With the Wach estimate W81T6 this proves the entire predicate
$\mathrm{TBat}$ on the stipulated level-one data for odd $p$.

In particular at the two places of H80M5 all kernel lattices have
the required $\mathrm{HypW}$, the middle terms are crystalline
on $G_1$, and $g_K=0$. The comparison of H80M3 lets one multiply
by the remaining nonnegative power to reach $c_B$. H80M4 supplies
the three rank-one clauses. Thus the original ramified
$\mathrm{admTB}$ is proved there for each good coefficient extension
and every tested row. The separate lifting estimate on a rank-$r$
row and the global lattice-index application have their own inputs.
% END K81C6
